% SGA 3, Expose V, Loop 1 English reconstruction.
% Scope: final diagram and conclusion of the affine constant-rank proof of
% Theorem 4.1, followed by Section 5(a) and the opening sentence of 5(b).
% Authority: Polo--Gille Exp5-13oct24.pdf, component-local p. 17
% (combined-reader p. 279), re-edited page header 267; internal markers
% 269--270. Jacob Reinhold's Markdown is comparison/drafting material only.

证明开始时所考虑的对偶图
\hyperref[fig:sga3-exp5-dual-012star]{\((0,1,2)^\ast\)} 导出如下交换图:

\SGAThreeDiagram[0.78\textwidth]
  {figures/exp5/p267_proof_iv_ring_modules_diagram.png}
  {完成\hyperref[thm:sga3-V-4.1-iv]{定理~4.1(iv)} 仿射情形证明的环与模的交换图.}
  {fig:sga3-exp5-p267-proof-iv-ring-modules}
  {Polo--Gille Exp5-13oct24.pdf, 分篇第 17 页;
  合订本第 279 页; 重版本页眉页码 267}

这里 \(u_0,u_1,u_2\) 分别具有分量对
\[
  (i,\varepsilon),\qquad
  (\delta_1,\delta_0\varepsilon),\qquad
  (\delta'_2,\delta'_0\delta_0\varepsilon).
\]
已知 \(u_1\) 是同构. 由于 \((0,1,2)^\ast\) 左边的两个方块是余笛卡儿图,
\(u_2\) 是双射. 图中的两个横行正合, 因此 \(u_0\) 是双射.%
\footnote{\label{note:sga3-V-26}编者注: 下面的论证是新增的.}
于是 \(A_0\) 是秩为 \(n\) 的局部自由 \(B\)-模.
由前面的各步归结可知,
\[
  \delta_0\otimes\delta_1:
  A_0\otimes_BA_0\xrightarrow{\ \sim\ }A_1
\]
是同构. 这就完成了所考虑的特殊情形下
\hyperref[thm:sga3-V-4.1]{定理~4.1} 的证明:
\(X_0\) 仿射, \(d_1\) 局部自由且秩恒为 \(n\).

\section*{5. 对有限平坦群胚取商 (一般情形)}
\phantomsection\label{sec:sga3-V-5}
\addcontentsline{toc}{section}{5. 对有限平坦群胚取商 (一般情形)}

\phantomsection\label{par:sga3-V-5-a}%
\noindent\textbf{a)}
设 \(U^{(n)}\) 是 \(X_0\) 中使 \(d_1\) 在其上有限局部自由且秩为 \(n\)
的最大开集. 概形 \(X_0\) 是各 \(U^{(n)}\) 的不交并. 两个笛卡儿方块

\begin{figure}[H]
  \centering
  \resizebox{0.29\textwidth}{!}{%
    \SGAThreeNativeInput{figures/exp5/p267_section5_rank_square_d0.png}}
  \qquad\textit{和}\qquad
  \resizebox{0.29\textwidth}{!}{%
    \SGAThreeNativeInput{figures/exp5/p267_section5_rank_square_d1.png}}
  \caption{控制秩分层 \(U^{(n)}\) 的两个笛卡儿方块.}
  \label{fig:sga3-exp5-p267-section5-rank-squares}
  \par\smallskip
  {\footnotesize\itshape 第 2 轮原生重构; 来源定位:
  Polo--Gille Exp5-13oct24.pdf, 分篇第 17 页;
  合订本第 279 页; 重版本页眉页码 267.}
\end{figure}

\noindent 表明 \(U^{(n)}\) 在 \(d_0\) 和 \(d_1\) 下的逆像
都等于 \(X_1\) 中使 \(d'_2\) 在其上局部自由且秩为 \(n\) 的最大开集.%
\footnote{\label{note:sga3-V-27}编者注: 由于 \(d_0\) (分别地, \(d_1\))
满射、平坦且有限, 因而忠实平坦且仿射,
\(d'_2\) 在一点 \(x\in X_1\) 的某邻域上秩为 \(n\),
当且仅当 \(d_1\) 在 \(d_0(x)\) (分别地, \(d_1(x)\))
的某邻域上秩为 \(n\).}
所以
\[
  d_0^{-1}(U^{(n)})=d_1^{-1}(U^{(n)}),
\]
且群胚 \(X_\ast\) 是由 \(X_\ast\) 在既开又闭子概形
\(U^{(n)}\) 上诱导的群胚 \(X_\ast^{(n)}\) 的不交并.
因此只须对每个 \(X_\ast^{(n)}\) 证明
\hyperref[thm:sga3-V-4.1]{定理~4.1}. 归结为 \(d_1\) 有限局部自由且秩为 \(n\) 的情形.

\medskip
\phantomsection\label{par:sga3-V-5-b}%
\noindent\textbf{b)}
\textit{现在可证明一般情形下的\hyperref[thm:sga3-V-4.1]{定理}.}

% Translation coverage: entire assigned component; adjoining components retain the original statement/proof seams.
% Source lines: 89; translated body lines before this audit comment: 80.
% Math preserved: 41 inline formulas, 3 display blocks, 0 equation environments; mathematical prose translated where present.
% Labels: 6/6; hyperrefs: 5/5; SGARefLinks: 0/0; inline targets: 0/0; footnotes: 2/2.
% Figure arguments preserved; all rendered captions and natural-language source locators translated.
% Unresolved TODOs: 0; compilation and visual inspection remain with the volume integrator.
